\begin{exercise}{2.5.2}
Consider the subset $\mathcal{D}$ of the set $\mathbb{N}\times\mathbb{N}$ of pairs of natural numbers defined by following axioms and rules
    \[(n,0) \in \mathcal{D}\]
    \[\frac{(n,n')\in \mathcal{D}}{(n, n+n')\in \mathcal{D}}\]
Use Rule Induction to prove
\[(n,n')\in \mathcal{D} \implies \exists k \in \mathbb{N} . n'=k\ast n\]
(where $\ast$ denotes multiplication). Use Mathematical Induction on $k$ to show conversely that if $n'=k\ast n$ then $(n,n') \in \mathcal{D}$

\end{exercise}

\begin{solution}
\begin{enumerate}
    \item \textbf{Rule induction} We claim \[(n,n')\in \mathcal{D} \implies \exists k \in \mathbb{N} . n'=k\ast n\]
We prove this using rule induction, which follows the way the set $\mathcal{D}$ is defined.
        \begin{enumerate}
            \item \textsc{Base case} The base rule says:
\[(n,0) \in \mathcal{D}\]
We need to show there exists 
$k\in \mathbb{N}$ such that $n'= k\ast n  = 0$. 
\\
Choose $k=0$. Then clearly:
\[0\ast n= 0 \]
Thus the base case holds.
            \item \textsc{Inductive step} The inductive rule says:
\[\frac{(n,n')\in \mathcal{D}}{(n, n+n')\in \mathcal{D}}\]

Assume as the induction hypothesis:
\[(n,n') \in \mathcal{D} \text{ and } n'=kn \implies n+n'= (k+1)\ast n\]
We must show: 
\[(n,n+n') \in \mathcal{D} \implies \exists k \in \mathbb{N}. n+n' = k\ast n\]
By induction hypothesis there exists some $k\in \mathbb{N}$ such that $n'=k\ast n$. Then:
\[n+n' = n+k\ast n = (k+1)\ast n\]
Let $k'=k+1$. Then $n+n' = k'\ast n$, with $k'\in \mathbb{N}$.
\\
Thus, the inductive step holds.
        
\end{enumerate}

By rule induction we can conclude 
\[(n,n')\in \mathcal{D} \implies \exists k \in \mathbb{N} . n'=k\ast n\]
\item \textbf{Mathematical Induction}
Now we prove the converse.
Claim:
\[\forall n \in \mathbb{N}, \forall k \in \mathbb{N}, (n, k\ast n) \in \mathcal{D}\]
We will use mathematical induction on $k$.
\begin{enumerate}
            \item \textsc{Base case $(k=0)$} We need to show:
\[(n,0\ast n)=(n,0)\in \mathcal{D}\]
This is exactly the base rule of $\mathcal{D}$, thus the base case holds.
            \item \textsc{Inductive step} 
            Assume the inductive hypothesis for some $k\in \mathbb{N}$: 
            \[(n, k\ast n) \in \mathcal{D}\]
            We need to prove: 
            \[(n, (k+1)\ast n) \in \mathcal{D}\]
            Observe that 
            \[(k+1) \ast n = k \ast n + n\]
            By inductive hypothesis $(n, k\ast n) \in \mathcal{D}$ and by inductive rule of $\mathcal{D}$, adding $n$ to the second component gives: 
            \[(n,k\ast n+n)=(n,(k+1)\ast n)\in\mathcal{D}\]
            Thus the inductive step holds.
\end{enumerate}
By mathematical induction on $k$ we conclude:

\[(n,k\ast n)\in\mathcal{D}  \qquad \forall k\in \mathbb{N},\forall n\in \mathbb{N}\]
\end{enumerate}
We conclude that $\mathcal{D}$ is exactly the set of pairs $(n,n')$ such that $n'=k\ast n$ for some $k\in\mathbb{N}$:
\[\mathcal{D}=\{(n,k\ast n)\mid n\in \mathbb{N},k\in \mathbb{N}\}\]
\end{solution}